% ===================================================================
%  Minimalbeispiel: Liedblatt mit Text und Akkorden – „Hänschen klein“
%  Text: Franz Wiedemann (1821–1882), Melodie: Volksweise (gemeinfrei), F-Dur.
%  Erste Strophe unter den Noten (musixlyr), Akkordsymbole über den Noten,
%  Strophen 2–3 zweispaltig darunter, alles auf einer Seite.
%  Eigenständig: braucht keine header.tex.
%
%  Drei Durchläufe:  pdflatex … → musixflx … → pdflatex …
%  (musixflx gleicht alle Notenzeilen auf dieselbe Länge aus.)
% ===================================================================
\documentclass[a4paper,12pt]{article}
\usepackage[T1]{fontenc}
\usepackage[utf8]{inputenc}
\usepackage[ngerman]{babel}
\usepackage{lmodern}
\usepackage[margin=1.8cm]{geometry}
\usepackage{multicol}       % Strophen zweispaltig
\usepackage{musixtex}
\input musixlyr             % Liedtext unter den Noten
\pagestyle{empty}
\setlength{\columnsep}{1em}

% Eine Strophe: \strophe{Nummer}{Text}; bleibt beim Spaltenumbruch zusammen.
\newcommand{\strophe}[2]{%
  \noindent\makebox[1.4em][l]{\textbf{#1.}}%
  \begin{minipage}[t]{\dimexpr\linewidth-1.4em}\raggedright #2\end{minipage}%
  \par\vspace{1.2em}}

\begin{document}
\begin{center}
{\LARGE\bfseries\sffamily Hänschen klein}\\[4pt]
{\small Text: Franz Wiedemann \textbullet\ Melodie: Volksweise \textbullet\ F-Dur}
\end{center}
\vspace{1em}

% Größeres Notensystem: Größenbefehl VOR \begin{music} (wirkt sonst nicht).
\Largemusicsize

\begin{music}
\parindent0pt                % erste Zeile nicht einrücken: alle Zeilen gleich lang
% Platz unter (Liedtext) und über (Akkorde) jedem System – INNERHALB von music:
\staffbotmarg=8\Interligne
\stafftopmarg=4\Interligne
\instrumentnumber{1}\setclef1\treble
\generalmeter{\meterfrac44}\generalsignature{-1}\nobarnumbers
\setsongraise1{8mm}         % Liedtext dicht unter die Hälse (bezogen auf den unteren Rand)
\setlyrics{strophe}{Häns-chen klein ging al-lein in die wei-te Welt hin-ein.
Stock und Hut steht ihm gut, ist gar wohl-ge-mut.
A-ber Mut-ter wei-net sehr, hat ja nun kein
Häns-chen mehr! »Wünsch dir Glück«, sagt ihr Blick, »kehr nur bald zu-rück!«}
\assignlyrics1{strophe}
\startpiece
% Tonhöhen als absolute Buchstaben: c = c', f g = f' g', h = a', i = h'/b',
% j = c''. Das b' ergibt sich aus der Tonart (\generalsignature{-1}).
% Akkorde mit \zcharnote{Höhe}{Symbol} nur an Harmoniewechseln.
\Notes\zcharnote{12}{\bf F}\qa j\qa h\ha h\en\bar          % Häns-chen klein
\Notes\zcharnote{12}{\bf C$^7$}\qa i\qa g\ha g\en\bar      % ging al-lein
\Notes\zcharnote{12}{\bf F}\qa f\qa g\qa h\qa i\en\bar     % in die wei-te
\Notes\qa j\qa j\ha j\en                                   % Welt hin-ein.
\alaligne
\Notes\zcharnote{12}{\bf F}\qa j\qa h\ha h\en\bar          % Stock und Hut
\Notes\zcharnote{12}{\bf C$^7$}\qa i\qa g\ha g\en\bar      % steht ihm gut,
\Notes\zcharnote{12}{\bf F}\qa f\qa h\qa j\qa j\en\bar     % ist gar wohl-ge-
\Notes\wh f\en                                             % mut.
\alaligne
\Notes\zcharnote{12}{\bf C$^7$}\qa g\qa g\qa g\qa g\en\bar % A-ber Mut-ter  
\Notes\qa g\qa h\ha i\en\bar                               % wei-net sehr,
\Notes\zcharnote{12}{\bf F}\qa h\qa h\qa h\qa h\en\bar     % hat ja nun kein
\Notes\qa h\qa i\ha j\en                                   % Häns-chen mehr!
\alaligne
\Notes\zcharnote{12}{\bf F}\qa j\qa h\ha h\en\bar          % »Wünsch dir Glück«,
\Notes\zcharnote{12}{\bf C$^7$}\qa i\qa g\ha g\en\bar      % sagt ihr Blick,
\Notes\zcharnote{12}{\bf F}\qa f\qa h\qa j\qa j\en\bar     % »kehr nur bald zu-
\Notes\wh f\en                                             % rück!«
\Endpiece                  % Schlussstrich
\end{music}
\assignlyrics1{}            % Zuweisung beenden (gilt sonst global weiter)

\vspace{2em}

\begin{multicols}{2}
\strophe{2}{Sieben Jahr, trüb und klar,\\ Hänschen in der Fremde war.\\
Da besinnt sich das Kind,\\ eilet heim geschwind.\\
Doch nun ist’s kein Hänschen mehr,\\ nein, ein großer Hans ist er,\\
braun gebrannt Stirn und Hand.\\ Wird er wohl erkannt?}

\columnbreak
\strophe{3}{Eins, zwei, drei gehn vorbei,\\ wissen nicht, wer das wohl sei.\\
Schwester spricht: »Welch Gesicht!«,\\ kennt den Bruder nicht.\\
Kommt daher die Mutter sein,\\ schaut ihm kaum ins Aug hinein,\\
ruft sie schon: »Hans, mein Sohn!\\ Grüß dich Gott, mein Sohn!«}
\end{multicols}

\end{document}
